\documentclass[UTF8,12pt]{ctexart}
\usepackage{amsmath,amssymb}
\usepackage{geometry}
\geometry{top=20mm,right=10mm,left=10mm,bottom=20mm}


\title{第一次作业报告}
\author{}
\date{}

\newcommand{\invn}{\tau}

\begin{document}
\maketitle

第一次作业在数的计算上问题不大, 但是在行列式的计算上经常忽略交换行或者列带来的负号, 以及在证明时语言表述不够清晰. 下面是对作业中出现问题较频繁的第 8、9、10、11 题的解答.
\section*{第 8 题}

\textbf{原题．}
设
\[
A=\begin{vmatrix}
a_{11}&a_{12}&a_{13}&a_{14}\\
a_{21}&a_{22}&a_{23}&a_{24}\\
a_{31}&a_{32}&a_{33}&a_{34}\\
a_{41}&a_{42}&a_{43}&a_{44}
\end{vmatrix}
\]
为四阶行列式, $A_{ij}\ (1\leq i,j\leq 4)$ 为 $(i,j)$ 元素的代数余子式. 写出下列算式为哪些行列式：
\begin{align*}
\text{(1)}\quad&
\sum_{j_1j_2j_3j_4}(-1)^{\invn(j_1j_2j_3j_4)}
a_{1j_1}b_{2j_2}a_{3j_3}a_{4j_4};\\
\text{(2)}\quad&
\sum_{j_1j_2j_3j_4}(-1)^{\invn(j_1j_2j_3j_4)}
a_{1j_1}a_{3j_2}a_{2j_3}a_{4j_4};\\
\text{(3)}\quad& a_{31}A_{21}+a_{32}A_{22}+a_{33}A_{23}+a_{34}A_{24};\\
\text{(4)}\quad& a_{13}A_{12}+a_{23}A_{22}+a_{33}A_{32}+a_{43}A_{42};\\
\text{(5)}\quad&
\sum_{j_1j_2j_3j_4}(-1)^{\invn(j_1j_2j_3j_4)}
a_{j_1 1}b_{j_2 2}a_{j_3 3}a_{j_4 4}.
\end{align*}
其中求和号表示对所有四阶排列 $j_1j_2j_3j_4$ 求和. 

\textbf{解．}

\textbf{(1)} 它对应于把第二行替换为$\{b_{2j_2}\}$
\[
\begin{vmatrix}
a_{11}&a_{12}&a_{13}&a_{14}\\
b_{21}&b_{22}&b_{23}&b_{24}\\
a_{31}&a_{32}&a_{33}&a_{34}\\
a_{41}&a_{42}&a_{43}&a_{44}
\end{vmatrix}.
\]

\textbf{(2)} 将 $A$ 的第二、三行互换, 所得行列式的定义展开式正是题给的和式. 因此
\[
\sum_{j_1j_2j_3j_4}(-1)^{\invn(j_1j_2j_3j_4)}
a_{1j_1}a_{3j_2}a_{2j_3}a_{4j_4}
=\begin{vmatrix}
a_{11}&a_{12}&a_{13}&a_{14}\\
a_{31}&a_{32}&a_{33}&a_{34}\\
a_{21}&a_{22}&a_{23}&a_{24}\\
a_{41}&a_{42}&a_{43}&a_{44}
\end{vmatrix}=-A.
\]

\textbf{(3)} 把 $A$ 的第二行换成第三行, 而其余各行不变. 沿第二行展开, 得到
\[
\begin{vmatrix}
a_{11}&a_{12}&a_{13}&a_{14}\\
a_{31}&a_{32}&a_{33}&a_{34}\\
a_{31}&a_{32}&a_{33}&a_{34}\\
a_{41}&a_{42}&a_{43}&a_{44}
\end{vmatrix}
=\sum_{j=1}^{4}a_{3j}A_{2j}
=a_{31}A_{21}+a_{32}A_{22}+a_{33}A_{23}+a_{34}A_{24}.
\]
这里各 $A_{2j}$ 未改变, 因为删去第二行后留下的三行与原行列式相同. 上述行列式有两行相同, 故结果为0.

\textbf{(4)} 把 $A$ 的第二列换成第三列, 而其余各列不变. 沿第二列展开, 得到
\[
\begin{vmatrix}
a_{11}&a_{13}&a_{13}&a_{14}\\
a_{21}&a_{23}&a_{23}&a_{24}\\
a_{31}&a_{33}&a_{33}&a_{34}\\
a_{41}&a_{43}&a_{43}&a_{44}
\end{vmatrix}
=\sum_{i=1}^{4}a_{i3}A_{i2}
=a_{13}A_{12}+a_{23}A_{22}+a_{33}A_{32}+a_{43}A_{42}.
\]
这里各 $A_{i2}$ 未改变, 因为删去第二列后留下的三列与原行列式相同. 上述行列式有两列相同, 故结果为0.

\textbf{(5)}  它对应于把第二列替换为$\{b_{j_2 2}\}$.
\[
\begin{vmatrix}
a_{11}&b_{21}&a_{31}&a_{41}\\
a_{12}&b_{22}&a_{32}&a_{42}\\
a_{13}&b_{23}&a_{33}&a_{43}\\
a_{14}&b_{24}&a_{34}&a_{44}
\end{vmatrix}.
\]

\section*{第 9 题}

\textbf{原题．} 按定义计算下列行列式：
\[
\text{(1)}\quad
D_1=\begin{vmatrix}
0&0&\cdots&0&1\\
0&0&\cdots&2&0\\
\vdots&\vdots&&\vdots&\vdots\\
0&n-1&\cdots&0&0\\
n&0&\cdots&0&0
\end{vmatrix};
\]
\[
\text{(2)}\quad
D_2=\begin{vmatrix}
0&1&0&\cdots&0\\
0&0&2&\cdots&0\\
\vdots&\vdots&\vdots&&\vdots\\
0&0&0&\cdots&n-1\\
n&0&0&\cdots&0
\end{vmatrix};
\]
\[
\text{(3)}\quad
D_3=\begin{vmatrix}
0&\cdots&0&1&0\\
0&\cdots&2&0&0\\
\vdots&&\vdots&\vdots&\vdots\\
n-1&\cdots&0&0&0\\
0&\cdots&0&0&n
\end{vmatrix}.
\]

\textbf{解．} 按行列式的定义, 每一项必须从每行、每列各取一个元素. 由于每一行都只有一个非零元素, 三个行列式的展开式都至多有一个非零项. 

\textbf{(1)} 唯一可能的非零项对应排列
\[
n,n-1,\ldots,2,1.
\]
该排列的逆序数为
\[
\invn(n,n-1,\ldots,1)=\binom n2=\frac{n(n-1)}2.
\]
因此
\[
D_1=(-1)^{\frac{n(n-1)}2}\,1\cdot2\cdots n
={(-1)^{\frac{n(n-1)}2}n!}.
\]

\textbf{(2)} 唯一可能的非零项对应排列
\[
2,3,\ldots,n,1.
\]
其中只有末尾的 $1$ 与前面的 $n-1$ 个数构成逆序, 故逆序数为 $n-1$. 因此
\[
D_2=(-1)^{n-1}\,1\cdot2\cdots n
={(-1)^{n-1}n!}.
\]

\textbf{(3)} 唯一可能的非零项对应排列
\[
n-1,n-2,\ldots,2,1,n.
\]
最后的 $n$ 不产生逆序；前 $n-1$ 个数恰为降序排列, 故
\[
\invn(n-1,n-2,\ldots,1,n)=\binom{n-1}{2}
=\frac{(n-1)(n-2)}2.
\]
因此
\[
D_3=(-1)^{\frac{(n-1)(n-2)}2}\,1\cdot2\cdots(n-1)\cdot n
={(-1)^{\frac{(n-1)(n-2)}2}n!}.
\]

\section*{第 10 题}

\textbf{原题．} 由行列式定义证明：
\[
D=\begin{vmatrix}
a_1&a_2&a_3&a_4&a_5\\
b_1&b_2&b_3&b_4&b_5\\
c_1&c_2&0&0&0\\
d_1&d_2&0&0&0\\
e_1&e_2&0&0&0
\end{vmatrix}=0.
\]

\textbf{证明．} 由五阶行列式的定义, 任一展开项均为
\[
(-1)^{\invn(j_1j_2j_3j_4j_5)}
a_{1j_1}a_{2j_2}a_{3j_3}a_{4j_4}a_{5j_5},
\]
其中 $j_1j_2j_3j_4j_5$ 是 $1,2,3,4,5$ 的一个排列（这里以 $a_{ij}$ 泛指矩阵的 $(i,j)$ 元素）. 

第三、四、五行的非零元素都只可能位于第一、二列. 因此, 若上述乘积非零, 就必须同时有
\[
j_3,j_4,j_5\in\{1,2\}.
\]
但排列中的 $j_3,j_4,j_5$ 必须彼此不同, 而集合 $\{1,2\}$ 只有两个元素, 不可能容纳三个互不相同的数. 故定义展开式中的每一项都为零, 从而$D=0$.

\section*{第 11 题}

\textbf{原题．} 由行列式定义计算
\[
f(x)=\begin{vmatrix}
2x&x&1&2\\
1&x&1&-1\\
3&2&x&1\\
1&1&1&x
\end{vmatrix}
\]
中 $x^4$ 与 $x^3$ 的系数, 并说明理由. 

\textbf{解．} 行列式定义展开式中的每一项, 都是从每行、每列各取一个元素所得的乘积（再乘以相应排列的符号）. 矩阵中含 $x$ 的位置为
\[
(1,1),(1,2),(2,2),(3,3),(4,4),
\]
其中 $(1,1)$ 处为 $2x$, 其余四处为 $x$. 

要产生四次项, 必须从四行、四列各取一个含 $x$ 的元素. 唯一可行的选取是主对角线
\[
(2x)\cdot x\cdot x\cdot x=2x^4,
\]
它对应恒等排列 $1234$, 符号为正. 因此$x^4$的系数为$2$.

要产生三次项, 需要从三个不同的行和列各取一个含 $x$ 的元素, 并从余下的一行、余下的一列取一个常数. 第三、四行若贡献 $x$, 就必须分别取第三、四列. 为了再取得两个不同列中的元素, 第一、二行只能分别取第二、第一列, 即唯一可行的选取为
\[
a_{12}a_{21}a_{33}a_{44}=x\cdot1\cdot x\cdot x=x^3.
\]
它对应排列 $2134$, 该排列有一个逆序, 符号为负. 因此$x^3$的系数为$-1$.

\end{document}
